\section{The Avellaneda--Stoikov Model}
\label{sec:model}

This section restates only what the empirical work needs. Full derivations are in
\cite{avellaneda2008} and \cite{amargos2026}.

\subsection{Setup}

Time runs over one trading session $t\in[0,T]$. We measure time in minutes, so $T=\Horizon$.
The mid-price is an arithmetic Brownian motion with no drift,
\begin{equation}
dS_t = \sigma\, dW_t ,
\label{eq:mid}
\end{equation}
where $\sigma$ is in rupees per $\sqrt{\text{minute}}$. The dealer posts a bid at $S_t-\delta^b_t$
and an ask at $S_t+\delta^a_t$, with distances $\delta^b,\delta^a\ge 0$. Market orders hit the
bid and lift the ask as two independent Poisson processes whose intensity depends only on the
distance of the quote from the mid,
\begin{equation}
\lambda(\delta) = A\,e^{-k\delta}.
\label{eq:lambda}
\end{equation}
$A$ is the fill rate of a quote sitting at the mid; $k$ is how fast that rate falls as the quote
moves away. Each fill is one share. Inventory $q_t$ is bids filled minus asks filled, cash $X_t$
moves by the fill prices, and the dealer maximises expected CARA utility of terminal wealth,
$\mathbb{E}\big[-e^{-\gamma(X_T+q_TS_T)}\big]$, with absolute risk aversion $\gamma$ in $1/\text{INR}$.

\subsection{The Optimal Quotes}

Solving the Hamilton--Jacobi--Bellman equation with the CARA ansatz gives two closed-form
objects~\cite{avellaneda2008}. The first is the \emph{reservation price}, the price at which the
dealer is indifferent to holding one more share,
\begin{equation}
r(s,q,t) = s - q\,\gamma\sigma^2 (T-t).
\label{eq:reservation}
\end{equation}
A long dealer ($q>0$) values the stock below the mid, a short dealer above it. The second is
the optimal \emph{total spread}, centred on $r$ rather than on $s$,
\begin{equation}
\delta^a+\delta^b = \underbrace{\gamma\sigma^2(T-t)}_{\text{inventory-risk premium}}
 \;+\; \underbrace{\frac{2}{\gamma}\ln\!\Big(1+\frac{\gamma}{k}\Big)}_{\text{market-making edge}} .
\label{eq:spread}
\end{equation}
The first term pays for carrying inventory risk over the time left and shrinks to zero at the
close. The second is what a dealer charges for liquidity alone; it depends only on how fast the
order book thins out ($k$). Figure~\ref{fig:schematic} draws both objects. Together they are the
whole strategy: quote $r\pm\frac12(\delta^a+\delta^b)$, floored so that neither quote crosses the
mid.

\begin{figure}[!htbp]
\centering
\begin{tikzpicture}[x=1cm,y=1cm,font=\small]
  \foreach \xo/\rr/\ttl in {0/2.9/{(a) Flat dealer, $q=0$}, 7.9/2.05/{(b) Long dealer, $q>0$}} {
    \begin{scope}[xshift=\xo cm]
      \node[anchor=west] at (-0.1,5.0) {\ttl};
      \draw[ink2, -{Stealth[length=2mm]}] (0,0.2) -- (0,4.6) node[above left, font=\footnotesize] {price};
      \draw[ink, line width=0.9pt] (0,2.9) -- (6.0,2.9) node[right, font=\footnotesize] {$s$};
      \draw[asdark, dashed, line width=0.8pt] (0,\rr) -- (6.0,\rr) node[right, font=\footnotesize] {$r$};
      \fill[asblue, opacity=0.13] (0.9,\rr-1.05) rectangle (4.4,\rr+1.05);
      \draw[asblue, line width=1.3pt] (0.9,\rr+1.05) -- (4.4,\rr+1.05) node[right, font=\footnotesize] {ask};
      \draw[asblue, line width=1.3pt] (0.9,\rr-1.05) -- (4.4,\rr-1.05) node[right, font=\footnotesize] {bid};
      \draw[decorate, decoration={brace, amplitude=4pt}, ink2] (0.75,2.9) -- (0.75,\rr+1.05)
        node[midway, left=4pt, font=\footnotesize] {$\delta^a$};
      \draw[decorate, decoration={brace, amplitude=4pt, mirror}, ink2] (0.75,2.9) -- (0.75,\rr-1.05)
        node[midway, left=4pt, font=\footnotesize] {$\delta^b$};
    \end{scope}
  }
  \draw[ink2, -{Stealth[length=2mm]}] (7.9+5.3,2.9) -- (7.9+5.3,2.05)
    node[midway, right, font=\footnotesize, align=left] {skew\\$q\gamma\sigma^2(T-t)$};
\end{tikzpicture}
\caption{How the optimal quotes move with inventory. (a) A flat dealer quotes symmetrically
around the mid. (b) A long dealer's reservation price falls below the mid, so the ask moves closer
(it sells more often) and the bid moves away (it buys less often). The width of the band is the
spread of equation~\eqref{eq:spread} and does not depend on $q$; only its position does.}
\label{fig:schematic}
\end{figure}

\subsection{Assumption Ledger}
\label{sec:ledger}

The base paper lists the assumptions behind these formulas. We restate the ones that real data can
test, numbered so that Section~\ref{sec:assumptions} can check each one.
\begin{enumerate}
\item[\textbf{A1}] \textbf{Exponential fill intensity:} $\lambda(\delta)=Ae^{-k\delta}$, equation~\eqref{eq:lambda}.
\item[\textbf{A2}] \textbf{Symmetric book:} the same $A$ and $k$ on the bid and the ask side.
\item[\textbf{A3}] \textbf{Driftless mid:} no predictable trend in $S_t$ over a session.
\item[\textbf{A4}] \textbf{Independent Gaussian increments:} $S_t$ is a Brownian motion, so its
  changes are uncorrelated, scale linearly in variance with the horizon, and are normal.
\item[\textbf{A5}] \textbf{Constant parameters:} $\sigma$, $A$ and $k$ do not change within
  a session or across sessions.
\item[\textbf{A6}] \textbf{No adverse selection:} a fill tells the dealer nothing about where the
  price goes next. This cannot be tested on its own; Section~\ref{sec:replay} measures its cost.
\end{enumerate}
